\begin{ex}
Tính kết quả của phép tính: $32 + 5$.
\shortans{37}
\loigiai{
Cộng hàng đơn vị: $2 + 5 = 7$. Giữ nguyên hàng chục là 3. Vậy $32 + 5 = 37$.
\textbf{Đáp số:} $37$.
}
\end{ex}

\begin{ex}
Tính kết quả của phép tính: $64 + 3$.
\shortans{67}
\loigiai{
Cộng hàng đơn vị: $4 + 3 = 7$. Giữ nguyên 6 chục. Vậy $64 + 3 = 67$.
\textbf{Đáp số:} $67$.
}
\end{ex}

\begin{ex}
Tính kết quả của phép tính: $81 + 7$.
\shortans{88}
\loigiai{
Cộng hàng đơn vị: $1 + 7 = 8$. Giữ nguyên 8 chục. Vậy $81 + 7 = 88$.
\textbf{Đáp số:} $88$.
}
\end{ex}

\begin{ex}
Điền số thích hợp vào ô trống: $43 + \dots = 48$.
\shortans{5}
\loigiai{
Muốn tìm số hạng chưa biết, lấy: $48 - 43 = 5$.
\textbf{Đáp số:} $5$.
}
\end{ex}

\begin{ex}
Điền số thích hợp vào ô trống: $\dots + 4 = 79$.
\shortans{75}
\loigiai{
Muốn tìm số hạng chưa biết, lấy: $79 - 4 = 75$.
\textbf{Đáp số:} $75$.
}
\end{ex}

\begin{ex}
Tính giá trị của biểu thức: $52 + 3 + 4$.
\shortans{59}
\loigiai{
Thực hiện: $52 + 3 = 55$; $55 + 4 = 59$.
\textbf{Đáp số:} $59$.
}
\end{ex}

\begin{ex}
Tính giá trị của biểu thức: $70 + 2 + 6$.
\shortans{78}
\loigiai{
Thực hiện: $70 + 2 = 72$; $72 + 6 = 78$.
\textbf{Đáp số:} $78$.
}
\end{ex}

\begin{ex}
Bạn Lan có 24 cái kẹo, mẹ cho thêm 5 cái kẹo nữa. Hỏi Lan có tất cả bao nhiêu cái kẹo?
\shortans{29}
\loigiai{
Lan có tất cả: $24 + 5 = 29$ (cái kẹo).
\textbf{Đáp số:} $29$.
}
\end{ex}

\begin{ex}
Lớp 1A có 31 học sinh, có thêm 4 học sinh mới chuyển đến. Hỏi lúc này lớp 1A có bao nhiêu học sinh?
\shortans{35}
\loigiai{
Số học sinh lớp 1A có lúc này là: $31 + 4 = 35$ (học sinh).
\textbf{Đáp số:} $35$.
}
\end{ex}

\begin{ex}
Bác nông dân nuôi 42 con gà, bác mua thêm 6 con gà nữa. Hỏi đàn gà nhà bác có tất cả bao nhiêu con gà?
\shortans{48}
\loigiai{
Đàn gà có tất cả: $42 + 6 = 48$ (con gà).
\textbf{Đáp số:} $48$.
}
\end{ex}

\begin{ex}
Tìm số tự nhiên $x$ thỏa mãn: $x - 3 = 65$.
\shortans{68}
\loigiai{
Ta có: $x = 65 + 3 = 68$.
\textbf{Đáp số:} $68$.
}
\end{ex}

\begin{ex}
Tìm số tự nhiên $y$ thỏa mãn: $y - 6 = 91$.
\shortans{97}
\loigiai{
Ta có: $y = 91 + 6 = 97$.
\textbf{Đáp số:} $97$.
}
\end{ex}

\begin{ex}
Vườn nhà Mai có 53 cây hoa hồng, bố trồng thêm 5 cây hoa hồng nữa. Hỏi trong vườn có tất cả bao nhiêu cây hoa hồng?
\shortans{58}
\loigiai{
Số cây hoa hồng là: $53 + 5 = 58$ (cây).
\textbf{Đáp số:} $58$.
}
\end{ex}

\begin{ex}
Một sợi dây dài 41 cm, nối thêm một đoạn dây dài 7 cm. Sợi dây mới dài bao nhiêu xăng-ti-mét?
\shortans{48}
\loigiai{
Chiều dài sợi dây mới là: $41 + 7 = 48$ (cm).
\textbf{Đáp số:} $48$.
}
\end{ex}

\begin{ex}
Tìm một số biết rằng lấy số đó trừ đi 4 thì được 52.
\shortans{56}
\loigiai{
Số cần tìm là: $52 + 4 = 56$.
\textbf{Đáp số:} $56$.
}
\end{ex}

\begin{ex}
Hai số có hiệu bằng 7, số bé là 82. Hỏi số lớn là số mấy?
\shortans{89}
\loigiai{
Số lớn là: $82 + 7 = 89$.
\textbf{Đáp số:} $89$.
}
\end{ex}

\begin{ex}
Số 45 cộng với số nào để được 49?
\shortans{4}
\loigiai{
Số đó là: $49 - 45 = 4$.
\textbf{Đáp số:} $4$.
}
\end{ex}

\begin{ex}
Số 71 cộng với số nào để được 77?
\shortans{6}
\loigiai{
Số đó là: $77 - 71 = 6$.
\textbf{Đáp số:} $6$.
}
\end{ex}

\begin{ex}
Một cửa hàng buổi sáng bán được 63 chiếc bút, buổi chiều bán được thêm 5 chiếc bút. Cả ngày cửa hàng bán được bao nhiêu chiếc bút?
\shortans{68}
\loigiai{
Cả ngày bán được: $63 + 5 = 68$ (chiếc bút).
\textbf{Đáp số:} $68$.
}
\end{ex}

\begin{ex}
Tính: $30 + 4 + 5$.
\shortans{39}
\loigiai{
Thực hiện: $30 + 4 = 34$; $34 + 5 = 39$.
\textbf{Đáp số:} $39$.
}
\end{ex}

\begin{ex}
Tìm số tự nhiên $a$ biết: $a + 2 = 86 + 1$.
\shortans{85}
\loigiai{
Vế phải: $86 + 1 = 87$. Ta có: $a + 2 = 87 \Rightarrow a = 87 - 2 = 85$.
\textbf{Đáp số:} $85$.
}
\end{ex}

\begin{ex}
So sánh: $51 + 6$ và $52 + 4$. Giá trị của biểu thức lớn hơn bằng bao nhiêu?
\shortans{57}
\loigiai{
$51 + 6 = 57$; $52 + 4 = 56$. Giá trị lớn hơn là $57$.
\textbf{Đáp số:} $57$.
}
\end{ex}

\begin{ex}
Hùng có 33 viên bi, Dũng cho Hùng thêm 6 viên bi. Sau đó Hùng làm mất 2 viên bi. Hỏi Hùng còn lại bao nhiêu viên bi?
\shortans{37}
\loigiai{
Sau khi nhận thêm Hùng có: $33 + 6 = 39$ (viên). Sau khi làm mất Hùng còn: $39 - 2 = 37$ (viên bi).
\textbf{Đáp số:} $37$.
}
\end{ex}

\begin{ex}
Băng giấy màu xanh dài 23 cm, cắt thêm 4 cm nữa để dài bằng băng giấy màu đỏ. Hỏi băng giấy màu đỏ dài bao nhiêu xăng-ti-mét?
\shortans{27}
\loigiai{
Băng giấy màu đỏ dài: $23 + 4 = 27$ (cm).
\textbf{Đáp số:} $27$.
}
\end{ex}

\begin{ex}
Tính: $92 + 1 + 2 + 3$.
\shortans{98}
\loigiai{
Thực hiện: $92 + 1 = 93$; $93 + 2 = 95$; $95 + 3 = 98$.
\textbf{Đáp số:} $98$.
}
\end{ex}
